% ======================================================================
% Chapter: The Lattice Construction — A Universal Engine for Examples
%
% This chapter establishes the systematic lattice construction framework
% that unifies all major vertex algebra examples in Part 2. Every family
% of chiral algebras studied in subsequent chapters — Heisenberg, affine
% Kac-Moody, W-algebras, deformation quantization — arises from lattice
% constructions or operations derived from them. The chapter provides
% the constructive machinery; the remaining chapters become applications.
% ======================================================================

\chapter{The Lattice Construction: A Universal Engine for Examples}
\label{ch:lattice}

The examples in Part~2 are not a miscellaneous collection. They arise from
a single source: \emph{lattice vertex algebras} and operations upon them.
Every vertex algebra family studied in the subsequent chapters --- Heisenberg,
affine Kac--Moody at level~1, $W$-algebras via quantum Drinfeld--Sokolov
reduction, deformation quantization via chiral Kontsevich formality --- is
either a lattice vertex algebra, a subalgebra of one, or obtained from one
by a functorial operation (extension, orbifold, BRST reduction).

The purpose of this chapter is to establish the lattice construction as a
\emph{universal engine}: a single framework from which the entire catalog
of examples can be derived systematically, with explicit presentations
amenable to bar complex computation and Koszul duality analysis.

\medskip
\noindent\textbf{Organization.}
Sections~\ref{sec:lattice:lattices}--\ref{sec:lattice:cocycles} recall
even lattices and their cocycle theory.
Section~\ref{sec:lattice:construction} constructs lattice vertex algebras
and classifies them by cocycle symmetry ($\Einf$ vs.\ $\Eone$).
Section~\ref{sec:lattice:frenkel-kac} proves the Frenkel--Kac--Segal
theorem identifying level-1 affine algebras with lattice vertex algebras.
Sections~\ref{sec:lattice:bar}--\ref{sec:lattice:koszul} compute bar
complexes and Koszul duals.
Section~\ref{sec:lattice:functorial} develops the functorial operations
(direct sum, extension, orbifold) that generate the full example catalog.
Sections~\ref{sec:lattice:hochschild}--\ref{sec:lattice:higher-genus}
treat chiral Hochschild cohomology and higher-genus aspects.


% ======================================================================
% Section 1: Even Lattices and Root Systems
% ======================================================================

\section{Even Lattices and Root Systems}
\label{sec:lattice:lattices}

\begin{definition}[Lattice]\label{def:lattice:lattice}
A \textbf{lattice} is a free abelian group $\Lambda$ of finite rank
equipped with a symmetric bilinear form
$\langle \cdot, \cdot \rangle \colon \Lambda \times \Lambda \to \Z$.
The lattice is:
\begin{enumerate}[label=(\roman*)]
\item \textbf{integral} if $\langle \alpha, \beta \rangle \in \Z$ for all
  $\alpha, \beta \in \Lambda$;
\item \textbf{even} if $\langle \alpha, \alpha \rangle \in 2\Z$ for all
  $\alpha \in \Lambda$ (this implies integrality);
\item \textbf{positive-definite} if $\langle \alpha, \alpha \rangle > 0$
  for all $\alpha \neq 0$;
\item \textbf{non-degenerate} if the map
  $\Lambda \to \Lambda^* := \Hom(\Lambda, \Z)$ given by
  $\alpha \mapsto \langle \alpha, \cdot \rangle$ is injective;
\item \textbf{unimodular} if this map is an isomorphism, i.e.,
  $\Lambda = \Lambda^*$.
\end{enumerate}
\end{definition}

\begin{definition}[Discriminant group]\label{def:lattice:discriminant}
For a non-degenerate lattice $\Lambda$, the \textbf{dual lattice} is
$\Lambda^* = \{\lambda \in \Lambda \otimes_\Z \bQ :
\langle \lambda, \alpha \rangle \in \Z \ \forall\, \alpha \in \Lambda\}$.
The \textbf{discriminant group} is the finite abelian group
$D(\Lambda) = \Lambda^*/\Lambda$,
whose order equals $|\det(\langle e_i, e_j \rangle)|$ for any basis
$\{e_i\}$ of $\Lambda$. The bilinear form on $\Lambda$ induces a
$\bQ/\Z$-valued bilinear form on $D(\Lambda)$.
\end{definition}

\begin{example}[Root lattices]\label{ex:lattice:root-lattices}
The ADE root lattices provide the fundamental examples.
\begin{center}
\begin{tabular}{l c c c c}
\textbf{Lattice} & \textbf{Rank} & $\det$ & $h^\vee$ & \textbf{Unimodular?} \\
\hline
$A_n$ & $n$ & $n+1$ & $n+1$ & No \\
$D_n$ & $n$ & $4$ & $2n-2$ & No \\
$E_6$ & $6$ & $3$ & $12$ & No \\
$E_7$ & $7$ & $2$ & $18$ & No \\
$E_8$ & $8$ & $1$ & $30$ & Yes
\end{tabular}
\end{center}
For each simply-laced root system $\Phi$, the root lattice
$\Lambda_\Phi = \Z\Phi$ is even and positive-definite.
The Cartan matrix $A_{ij} = \langle \alpha_i, \alpha_j \rangle$ encodes
the bilinear form on simple roots.
\end{example}

\begin{remark}[Uniqueness of $E_8$]\label{rem:lattice:e8-unique}
By Mordell's theorem (1938), $E_8$ is the unique even unimodular
positive-definite lattice of rank~8. In rank~16, there are exactly two
such lattices: $E_8 \oplus E_8$ and $D_{16}^+$ (the latter being the
overlattice of $D_{16}$ by its spinor class). This classification is
the lattice-theoretic origin of the two heterotic string theories.
\end{remark}

\begin{definition}[Lattice operations]\label{def:lattice:operations}
We record the operations that will become functorial on vertex algebras
in Section~\ref{sec:lattice:functorial}.
\begin{enumerate}[label=(\roman*)]
\item \textbf{Direct sum.}
  For lattices $\Lambda_1, \Lambda_2$, the orthogonal direct sum
  $\Lambda_1 \oplus \Lambda_2$ has
  $\langle (\alpha_1, \alpha_2), (\beta_1, \beta_2) \rangle =
  \langle \alpha_1, \beta_1 \rangle + \langle \alpha_2, \beta_2 \rangle$.
\item \textbf{Sublattice.}
  A sublattice $\Lambda' \hookrightarrow \Lambda$ of finite index
  gives rise to $D(\Lambda') \twoheadrightarrow D(\Lambda)$.
\item \textbf{Overlattice (gluing).}
  For an isotropic subgroup $I \subset D(\Lambda)$ (meaning
  $\langle \gamma, \gamma' \rangle \equiv 0 \bmod \Z$ for all
  $\gamma, \gamma' \in I$), the preimage
  $\Lambda_I = \{\lambda \in \Lambda^* : [\lambda] \in I\}$ is a
  lattice with $\Lambda \subset \Lambda_I \subset \Lambda^*$.
\item \textbf{Orbifold.}
  For $\sigma \in \operatorname{Aut}(\Lambda)$, the fixed sublattice
  $\Lambda^\sigma$ and the coinvariants $\Lambda/(1-\sigma)\Lambda$
  both carry induced bilinear forms.
\end{enumerate}
\end{definition}


% ======================================================================
% Section 2: Cocycle Theory
% ======================================================================

\section{Cocycle Theory}
\label{sec:lattice:cocycles}

The vertex algebra structure on $V_\Lambda$ requires a choice of
2-cocycle on $\Lambda$. Different cocycles produce genuinely different
algebraic structures, ranging from honest vertex algebras ($\Einf$-chiral)
to nonlocal $\Eone$-chiral algebras. This section develops the cocycle
theory that controls this distinction.

\begin{definition}[2-cocycle on a lattice]\label{def:lattice:2-cocycle}
A \textbf{2-cocycle} on a lattice $\Lambda$ is a function
$\varepsilon \colon \Lambda \times \Lambda \to \C^\times$ satisfying:
\begin{equation}\label{eq:lattice:cocycle-condition}
\varepsilon(\alpha, \beta)\,\varepsilon(\alpha + \beta, \gamma)
= \varepsilon(\alpha, \beta + \gamma)\,\varepsilon(\beta, \gamma)
\end{equation}
for all $\alpha, \beta, \gamma \in \Lambda$.
The cocycle is:
\begin{enumerate}[label=(\roman*)]
\item \textbf{normalized} if
  $\varepsilon(\alpha, 0) = \varepsilon(0, \alpha) = 1$ for all $\alpha$;
\item \textbf{bimultiplicative} if
  $\varepsilon(\alpha + \alpha', \beta) =
  \varepsilon(\alpha, \beta)\,\varepsilon(\alpha', \beta)$
  and similarly in the second argument;
\item \textbf{symmetric} (in the sense of Borcherds) if
\begin{equation}\label{eq:lattice:borcherds-symmetry}
\varepsilon(\alpha, \beta)
= (-1)^{\langle \alpha, \beta \rangle}\,\varepsilon(\beta, \alpha)
\end{equation}
for all $\alpha, \beta \in \Lambda$.
\end{enumerate}
\end{definition}

\begin{remark}\label{rem:lattice:symmetry-convention}
The symmetry condition~\eqref{eq:lattice:borcherds-symmetry} is the
\emph{Borcherds axiom}: it is precisely the condition needed for the
vertex operator skew-symmetry $Y(a,z)b = e^{z\partial}Y(b,-z)a$ to hold.
Note that for an even lattice, $(-1)^{\langle\alpha,\beta\rangle}$
is well-defined. For an odd lattice, one would need
$(-1)^{\langle\alpha,\alpha\rangle\langle\beta,\beta\rangle/4}$
corrections; we restrict to even lattices throughout.
\end{remark}

\begin{definition}[Commutator map]\label{def:lattice:commutator}
The \textbf{commutator} of a cocycle $\varepsilon$ is
\[
c_\varepsilon(\alpha, \beta)
:= \frac{\varepsilon(\alpha, \beta)}{\varepsilon(\beta, \alpha)}.
\]
\end{definition}

\begin{lemma}[Cocycle classification; \ClaimStatusProvedHere]\label{lem:lattice:cocycle-class}
Let $\Lambda$ be an even lattice. Then:
\begin{enumerate}[label=(\roman*)]
\item For any normalized bimultiplicative 2-cocycle $\varepsilon$, the
  commutator satisfies
  $c_\varepsilon(\alpha, \beta) = (-1)^{\langle\alpha,\beta\rangle}$
  if and only if $\varepsilon$ is symmetric.
\item The set of symmetric cocycles on $\Lambda$ is a torsor for
  $\Hom(\Lambda \otimes \Lambda, \C^\times)^{\text{sym}}$, the group of
  symmetric bicharacters.
\item The cohomology group $H^2(\Lambda, \C^\times)$ classifies
  cocycles up to coboundary equivalence. Two cocycles
  $\varepsilon, \varepsilon'$ are cohomologous if there exists
  $f \colon \Lambda \to \C^\times$ with
  $\varepsilon'(\alpha,\beta) =
  f(\alpha)\,f(\beta)\,f(\alpha+\beta)^{-1}\,\varepsilon(\alpha,\beta)$.
\end{enumerate}
\end{lemma}

\begin{proof}
(i) Computing from the Borcherds symmetry:
$c_\varepsilon(\alpha,\beta) =
\varepsilon(\alpha,\beta)/\varepsilon(\beta,\alpha) =
(-1)^{\langle\alpha,\beta\rangle}$.
Conversely, if $c_\varepsilon(\alpha,\beta) =
(-1)^{\langle\alpha,\beta\rangle}$, then
$\varepsilon(\alpha,\beta) =
(-1)^{\langle\alpha,\beta\rangle}\varepsilon(\beta,\alpha)$
is the Borcherds axiom.

(ii) If $\varepsilon, \varepsilon'$ are both symmetric, then
$\varepsilon'/\varepsilon$ is a symmetric bicharacter. Conversely,
multiplying a symmetric cocycle by a symmetric bicharacter preserves
the Borcherds axiom.

(iii) Standard group cohomology: $\Lambda \cong \Z^d$ is free abelian,
so $H^2(\Lambda, \C^\times) \cong (\C^\times)^{\binom{d}{2}}$,
parametrized by the values $\varepsilon(e_i, e_j)$ for $i < j$ on a
basis $\{e_1, \ldots, e_d\}$.
\end{proof}

\begin{construction}[Standard symmetric cocycle]\label{constr:lattice:standard-cocycle}
For an even lattice $\Lambda$ with ordered basis $\{e_1, \ldots, e_d\}$,
the \textbf{standard symmetric cocycle} is:
\begin{equation}\label{eq:lattice:standard-cocycle}
\varepsilon_0(\alpha, \beta)
= (-1)^{\sum_{i < j} \alpha_i \beta_j}
\end{equation}
where $\alpha = \sum_i \alpha_i e_i$ and $\beta = \sum_j \beta_j e_j$.
One verifies:
\[
\varepsilon_0(\alpha,\beta)\,\varepsilon_0(\beta,\alpha)
= (-1)^{\sum_{i<j}(\alpha_i\beta_j + \beta_i\alpha_j)}
= (-1)^{\langle\alpha,\beta\rangle - \sum_i \alpha_i\beta_i\langle e_i,e_i\rangle/1},
\]
which for an even lattice reduces to
$(-1)^{\langle\alpha,\beta\rangle}$, confirming the Borcherds axiom.
\end{construction}

\begin{example}[Non-symmetric cocycle on $\Z^2$]\label{ex:lattice:nonsym-z2}
On $\Lambda = \Z^2$ with standard bilinear form
$\langle (a,b), (c,d) \rangle = ac + bd$, define:
\[
\varepsilon((a,b), (c,d)) = e^{2\pi i \theta \cdot ad}
\]
for $\theta \in \R \setminus \bQ$.
The cocycle condition holds since
$\varepsilon((a,b),(c,d))\,\varepsilon((a+c,b+d),(e,f))
= e^{2\pi i\theta(ad + (a+c)f)}$
equals
$\varepsilon((a,b),(c+e,d+f))\,\varepsilon((c,d),(e,f))
= e^{2\pi i\theta(a(d+f) + cf)}$.

For the symmetry check:
\[
\frac{\varepsilon((a,b),(c,d))}{\varepsilon((c,d),(a,b))}
= e^{2\pi i\theta(ad - cb)},
\]
which is not $(-1)^{ac+bd}$ for irrational $\theta$. Hence this cocycle
is not symmetric, producing a strictly $\Eone$-chiral algebra
(see Theorem~\ref{thm:lattice:e1-vs-einf}).
\end{example}

\begin{remark}[Non-symmetric deformations via roots of unity]
\label{rem:lattice:root-of-unity-deformation}
More generally, given a symmetric cocycle $\varepsilon_0$ and an
antisymmetric bilinear form $q \colon \Lambda \times \Lambda \to \Z/N\Z$,
the product
$\varepsilon(\alpha, \beta) = \varepsilon_0(\alpha, \beta) \cdot
\zeta^{q(\alpha,\beta)}$
(where $\zeta = e^{2\pi i/N}$) is a non-symmetric cocycle whenever $q$
is not identically zero. Such deformations interpolate between the
commutative ($\Einf$) and noncommutative ($\Eone$) worlds.
\end{remark}


% ======================================================================
% Section 3: Lattice Vertex Algebra Construction
% ======================================================================

\section{Lattice Vertex Algebra Construction}
\label{sec:lattice:construction}

\begin{construction}[Lattice vertex algebra]\label{constr:lattice:va}
Let $\Lambda$ be an even positive-definite lattice of rank~$d$ with
bilinear form $\langle \cdot, \cdot \rangle$, and let
$\varepsilon \colon \Lambda \times \Lambda \to \C^\times$ be a
normalized 2-cocycle. The \textbf{lattice vertex algebra}
$\Vlat_\Lambda^\varepsilon$ is defined as follows.

\textbf{State space.}
Let $\cH = \C[\alpha^i_{-n} : 1 \leq i \leq d,\; n \geq 1]$ be the
Heisenberg Fock space (bosonic oscillator algebra) for
$\Lambda \otimes_\Z \C$, with vacuum $|0\rangle$ annihilated by all
$\alpha^i_n$ for $n > 0$. Let $\C[\Lambda]$ denote the group algebra
with basis $\{e^\alpha : \alpha \in \Lambda\}$ and multiplication
$e^\alpha \cdot e^\beta = \varepsilon(\alpha,\beta)\, e^{\alpha+\beta}$.
The state space is:
\begin{equation}\label{eq:lattice:state-space}
\Vlat_\Lambda^\varepsilon
= \cH \otimes_\C \C[\Lambda]
= \bigoplus_{\alpha \in \Lambda} \cH \otimes e^\alpha.
\end{equation}
This is $\Lambda$-graded by declaring
$\deg(\cH \otimes e^\alpha) = \alpha$.

\textbf{Vertex operators.}
For $\alpha \in \Lambda$, the vertex operator is:
\begin{equation}\label{eq:lattice:vertex-op}
Y(e^\alpha, z) = E^-(\alpha, z)\, E^+(\alpha, z)\, e^\alpha\, z^{\alpha_0}
\end{equation}
where:
\begin{align*}
E^-(\alpha, z)
&= \exp\!\Bigl(-\sum_{n < 0} \frac{\alpha_n}{n}\, z^{-n}\Bigr)
  \quad \text{(creation part)}, \\
E^+(\alpha, z)
&= \exp\!\Bigl(-\sum_{n > 0} \frac{\alpha_n}{n}\, z^{-n}\Bigr)
  \quad \text{(annihilation part)},
\end{align*}
the operator $e^\alpha$ acts on $\C[\Lambda]$ by
$e^\alpha \cdot e^\beta = \varepsilon(\alpha,\beta)\,e^{\alpha+\beta}$,
and $z^{\alpha_0}$ is the formal variable tracking the zero mode.

\textbf{OPE.}
\begin{equation}\label{eq:lattice:ope}
Y(e^\alpha, z)\, Y(e^\beta, w)
= \varepsilon(\alpha, \beta)\, (z-w)^{\langle\alpha,\beta\rangle}\,
  Y(e^{\alpha+\beta}, w) + \text{regular terms}.
\end{equation}

\textbf{Conformal vector.}
Choose an orthonormal basis $\{h_i\}$ of $\Lambda \otimes \C$.
The conformal vector is
$\omega = \frac{1}{2}\sum_{i=1}^d h_{i,{-1}}^2 |0\rangle$,
giving rise to the Virasoro algebra with central charge $c = d = \operatorname{rank}(\Lambda)$.
The conformal weight of $e^\alpha$ is:
\begin{equation}\label{eq:lattice:conformal-weight}
h_{e^\alpha} = \frac{\langle\alpha,\alpha\rangle}{2}.
\end{equation}
\end{construction}

\begin{theorem}[$\Eone$ vs.\ $\Einf$ classification; \ClaimStatusProvedHere]
\label{thm:lattice:e1-vs-einf}
The lattice vertex algebra $\Vlat_\Lambda^\varepsilon$ is:
\begin{enumerate}[label=(\roman*)]
\item an $\Einf$-chiral algebra (a vertex algebra in the sense of
  Borcherds) if and only if $\varepsilon$ is symmetric;
\item an $\Eone$-chiral algebra for any cocycle $\varepsilon$;
\item strictly $\Eone$ (not $\Einf$) when $\varepsilon$ is not symmetric.
\end{enumerate}
\end{theorem}

\begin{proof}
\textbf{(ii)} The cocycle condition~\eqref{eq:lattice:cocycle-condition}
on $\varepsilon$ ensures associativity of the vertex operator products
$Y(Y(e^\alpha,z-w)\,e^\beta, w) = Y(e^\alpha, z)\,Y(e^\beta, w)$.
Hence $\Vlat_\Lambda^\varepsilon$ is always at least $\Eone$-chiral.

\textbf{(i)} and \textbf{(iii)}: The skew-symmetry axiom for vertex
algebras requires
$Y(e^\alpha, z)\, e^\beta = e^{z\partial} Y(e^\beta, -z)\, e^\alpha$.
Computing both sides:
\begin{align*}
\text{LHS}
&= \varepsilon(\alpha,\beta)\, z^{\langle\alpha,\beta\rangle}\,
  e^{\alpha+\beta} + \cdots, \\
\text{RHS}
&= \varepsilon(\beta,\alpha)\, (-z)^{\langle\beta,\alpha\rangle}\,
  e^{\alpha+\beta} + \cdots
= \varepsilon(\beta,\alpha)\,(-1)^{\langle\alpha,\beta\rangle}\,
  z^{\langle\alpha,\beta\rangle}\, e^{\alpha+\beta} + \cdots.
\end{align*}
Equality of the leading terms requires
$\varepsilon(\alpha,\beta) =
(-1)^{\langle\alpha,\beta\rangle}\,\varepsilon(\beta,\alpha)$,
which is the Borcherds symmetry condition~\eqref{eq:lattice:borcherds-symmetry}.
Hence skew-symmetry holds if and only if $\varepsilon$ is symmetric.
When $\varepsilon$ is not symmetric, the OPE exchange picks up a phase
$c_\varepsilon(\alpha,\beta)/(-1)^{\langle\alpha,\beta\rangle} \neq 1$,
so the algebra is strictly $\Eone$.
\end{proof}

\begin{example}[Heisenberg as a lattice vertex algebra]
\label{ex:lattice:heisenberg}
The Heisenberg vertex algebra $\cH_\kappa$ at level $\kappa$ is the
$\alpha = 0$ sector of $\Vlat_\Lambda^\varepsilon$ for
\emph{any} lattice $\Lambda$:
\[
\cH_d := \Vlat_0 = \cH \otimes e^0 \cong \C[\alpha^i_{-n} : i, n \geq 1],
\]
with $d = \operatorname{rank}(\Lambda)$ free bosons. The cocycle
$\varepsilon$ does not enter the $\alpha = 0$ sector. This is the
simplest output of the lattice engine, and the base case for all
bar complex computations.
\end{example}


% ======================================================================
% Section 4: The Frenkel-Kac-Segal Theorem
% ======================================================================

\section{The Frenkel--Kac--Segal Theorem}
\label{sec:lattice:frenkel-kac}

The deepest structural result in the theory of lattice vertex algebras
is the theorem of Frenkel--Kac \cite{FK80} and Segal \cite{Se81},
identifying level-1 representations of simply-laced affine Lie algebras
with lattice vertex algebras on root lattices.

\begin{theorem}[Frenkel--Kac--Segal; \ClaimStatusProvedElsewhere]\label{thm:lattice:frenkel-kac}
Let $\fg$ be a simply-laced simple Lie algebra with root lattice
$\Lambda_\fg$. Let $\varepsilon$ be a symmetric cocycle on $\Lambda_\fg$
satisfying $\varepsilon(\alpha,\beta)\,\varepsilon(\beta,\alpha) =
(-1)^{\langle\alpha,\beta\rangle}$. Then there is an isomorphism of
vertex algebras:
\begin{equation}\label{eq:lattice:frenkel-kac}
\Vlat_{\Lambda_\fg}^\varepsilon
\;\cong\; L(\AffKM{g}, 1)
\end{equation}
where $L(\AffKM{g}, 1)$ denotes the simple quotient of the level-1
vacuum module for the affine Kac--Moody algebra $\AffKM{g}$.
\end{theorem}

\begin{proof}[Proof sketch]
The proof proceeds in three steps.

\textbf{Step 1: Embedding $\AffKM{g}$ into $\Vlat_{\Lambda_\fg}$.}
The Heisenberg subalgebra $\cH \subset \Vlat_{\Lambda_\fg}$ provides
the Cartan subalgebra $\AffKM{h} \hookrightarrow \AffKM{g}$.
For each root $\alpha \in \Phi$, the vertex operator $Y(e^\alpha, z)$
provides the root current. Specifically, define:
\[
E_\alpha(z) := Y(e^\alpha, z), \quad
H_i(z) := Y(\alpha_{i,-1}|0\rangle, z)
\]
for simple roots $\alpha_i$.

\textbf{Step 2: Verifying the OPE.}
For simple roots $\alpha_i, \alpha_j$ with
$\langle \alpha_i, \alpha_j \rangle = A_{ij}$ (Cartan matrix entry),
the lattice OPE~\eqref{eq:lattice:ope} gives:
\[
E_{\alpha_i}(z)\, E_{-\alpha_i}(w)
= \frac{1}{(z-w)^2} + \frac{H_i(w)}{z-w} + \text{regular},
\]
reproducing the $\AffKM{g}$ current algebra at level $k = 1$
(the level equals $\langle\alpha,\alpha\rangle/2 = 1$ for
simply-laced roots).

\textbf{Step 3: Simplicity.}
The lattice vertex algebra $\Vlat_{\Lambda_\fg}$ is simple (it has no
proper graded ideals, since $\Lambda_\fg$ is the root lattice of a
simple Lie algebra). Hence the image of the embedding
$L(\AffKM{g},1) \hookrightarrow \Vlat_{\Lambda_\fg}$ is the
entire algebra.
\end{proof}

\begin{example}[Explicit verification for $\AffKM{sl}_2$ at level 1]
\label{ex:lattice:sl2-fks}
The $A_1$ root lattice is $\Lambda = \Z\alpha$ with
$\langle\alpha,\alpha\rangle = 2$.
The lattice vertex algebra $\Vlat_\Lambda$ has generators:
\begin{itemize}
\item $h(z) := Y(\alpha_{-1}|0\rangle, z)$ --- the Cartan current,
\item $e(z) := Y(e^\alpha, z)$ --- the raising current,
\item $f(z) := Y(e^{-\alpha}, z)$ --- the lowering current.
\end{itemize}
The OPE from~\eqref{eq:lattice:ope}:
\begin{alignat*}{2}
h(z)\,h(w) &= \frac{2}{(z-w)^2} + \text{reg}
  &\qquad &\text{(level $k = 1$, $\langle\alpha,\alpha\rangle = 2$)}, \\
e(z)\,f(w)
&= \frac{\varepsilon(\alpha,-\alpha)}{(z-w)^2}
  + \frac{h(w)}{z-w} + \text{reg}
  &\qquad &\text{(simple pole gives $h$)}, \\
h(z)\,e(w) &= \frac{2\,e(w)}{z-w} + \text{reg}
  &\qquad &\text{(root $\alpha$ has $\langle\alpha,\alpha\rangle=2$)}.
\end{alignat*}
With the normalization $\varepsilon(\alpha,-\alpha) = 1$ (which is
achievable by coboundary adjustment), this reproduces the $\AffKM{sl}_2$
OPE at level $k=1$. The Sugawara construction gives
$T = \frac{1}{2(k+h^\vee)}(\frac{1}{2}\normord{h\,h} +
\normord{e\,f} + \normord{f\,e})$
with $h^\vee = 2$, central charge $c = 3 \cdot 1/(1+2) = 1$,
matching $\operatorname{rank}(\Lambda) = 1$.
\end{example}

\begin{remark}[Non-simply-laced algebras]\label{rem:lattice:non-simply-laced}
For non-simply-laced $\fg$ (types $B_n$, $C_n$, $F_4$, $G_2$),
the Frenkel--Kac--Segal construction does not apply directly, because
the root lattice contains roots of different lengths with
$\langle\alpha,\alpha\rangle \neq 2$. Instead, one obtains
non-simply-laced affine algebras via \emph{orbifold} constructions:
start with a simply-laced lattice $\Lambda$ and quotient by a
diagram automorphism $\sigma$. For instance:
\begin{itemize}
\item $\AffKM{sp}_{2n}$ at level 1 arises from $D_{n+1}$ modulo its
  order-2 diagram automorphism.
\item $\widehat{G}_2$ at level 1 arises from $D_4$ modulo its order-3
  triality automorphism.
\end{itemize}
See Section~\ref{sec:lattice:functorial} for the orbifold construction.
\end{remark}


% ======================================================================
% Section 5: Bar Complexes of Lattice Vertex Algebras
% ======================================================================

\section{Bar Complexes of Lattice Vertex Algebras}
\label{sec:lattice:bar}

We now apply the geometric bar construction of Part~1 to lattice vertex
algebras. The lattice grading provides a powerful organizational tool:
the bar differential preserves the total lattice degree, decomposing the
bar complex into orthogonal summands indexed by $\Lambda$.

\begin{construction}[Bar complex of $\Vlat_\Lambda^\varepsilon$]
\label{constr:lattice:bar-complex}
The geometric bar complex $\barBgeom(\Vlat_\Lambda^\varepsilon)$ has:

\textbf{Generators.} Elements of the form
\[
[e^{\alpha_1} \mid e^{\alpha_2} \mid \cdots \mid e^{\alpha_p}]
\otimes \omega
\]
where $\alpha_i \in \Lambda \setminus \{0\}$ and $\omega$ is a
logarithmic form on $\overline{C}_p(X)$. (We suppress Heisenberg
descendants $\alpha_{-n}^i |0\rangle$ for notational clarity; they
contribute additional tensor factors.)

\textbf{Grading.} The total lattice degree is $\sum_{i=1}^p \alpha_i \in \Lambda$, and the bar degree is $p$.

\textbf{Differential.}
The bar differential $d = d_{\text{res}}$ sums residues over collision
divisors. For a degree-2 element $[e^\alpha \mid e^\beta] \otimes 1$:
\begin{equation}\label{eq:lattice:bar-diff}
d([e^\alpha \mid e^\beta] \otimes 1)
= \varepsilon(\alpha,\beta)\,[e^{\alpha+\beta}] \otimes 1
\end{equation}
when $\langle\alpha,\beta\rangle = -1$ (simple pole in the OPE), and
$d = 0$ when $\langle\alpha,\beta\rangle \geq 0$ (no pole) or
$\langle\alpha,\beta\rangle \leq -2$ (higher poles pair trivially
with $d\log$).
\end{construction}

\begin{theorem}[Lattice bar complex structure; \ClaimStatusProvedHere]
\label{thm:lattice:bar-structure}
The bar complex $\barBgeom(\Vlat_\Lambda^\varepsilon)$ satisfies:
\begin{enumerate}[label=(\roman*)]
\item \textbf{Lattice grading.}
  The differential preserves the total lattice degree
  $\sum_i \alpha_i \in \Lambda$, decomposing the bar complex as
  $\barBgeom(\Vlat_\Lambda^\varepsilon) =
  \bigoplus_{\gamma \in \Lambda}
  \barBgeom(\Vlat_\Lambda^\varepsilon)_\gamma$.
\item \textbf{Residue criterion.}
  Simple poles occur only when $\langle\alpha_i,\alpha_j\rangle = -1$
  for some pair, and these are the only contributions to the differential.
\item \textbf{Cocycle dependence.}
  The cocycle $\varepsilon$ appears explicitly in the bar differential
  via~\eqref{eq:lattice:bar-diff}. Different cocycles give
  non-isomorphic bar complexes.
\end{enumerate}
\end{theorem}

\begin{proof}
(i) The OPE $Y(e^\alpha, z)\,Y(e^\beta, w)$ produces vertex operators
of lattice degree $\alpha + \beta$. The residue operation preserves
this degree.

(ii) The OPE~\eqref{eq:lattice:ope} has a pole of order
$-\langle\alpha,\beta\rangle$ at $z = w$. Pairing with the logarithmic
form $\eta_{ij} = d\log(z_i - z_j)$ extracts the residue, which is
nonzero only for a simple pole (order 1, i.e.,
$\langle\alpha,\beta\rangle = -1$).

(iii) The factor $\varepsilon(\alpha,\beta)$ in the
OPE~\eqref{eq:lattice:ope} passes through to the residue
formula~\eqref{eq:lattice:bar-diff}, so the differential depends on
the cocycle.
\end{proof}

\begin{computation}[$A_1$ bar complex]\label{comp:lattice:bar-A1}
For $\Lambda = A_1 = \Z\alpha$ with $\langle\alpha,\alpha\rangle = 2$,
the relevant inner products are:
\[
\langle n\alpha, m\alpha \rangle = 2nm.
\]
Since $\langle n\alpha, m\alpha\rangle = -1$ has no integer solutions,
\emph{no simple poles occur}. The bar differential vanishes identically,
and $H_*(\barBgeom(\Vlat_{A_1})) \cong \barBgeom(\Vlat_{A_1})$ as
graded vector spaces.

This reflects the fact that the $\AffKM{sl}_2$ current algebra at
level~1 is ``too small'' to have nontrivial collisions in bar degree~2
within a single lattice sector. The nontrivial bar differential involves
Heisenberg descendants (which carry additional oscillator modes that can
produce simple poles via derivative OPEs).
\end{computation}

\begin{computation}[$A_2$ bar complex]\label{comp:lattice:bar-A2}
For $\Lambda = A_2$ with simple roots $\alpha_1, \alpha_2$ satisfying
$\langle\alpha_1,\alpha_1\rangle = \langle\alpha_2,\alpha_2\rangle = 2$
and $\langle\alpha_1,\alpha_2\rangle = -1$:

\textbf{Bar degree 1:} Generated by $[e^\alpha]$ for
$\alpha \in \Lambda \setminus \{0\}$.

\textbf{Bar degree 2:} The key differential is
\[
d([e^{\alpha_1} \mid e^{\alpha_2}] \otimes 1)
= \varepsilon(\alpha_1,\alpha_2)\,[e^{\alpha_1+\alpha_2}] \otimes 1
\]
since $\langle\alpha_1,\alpha_2\rangle = -1$. This maps the element
$[e^{\alpha_1} \mid e^{\alpha_2}]$ of lattice degree
$\alpha_1 + \alpha_2$ to $[e^{\alpha_1+\alpha_2}]$ of the same lattice
degree, with coefficient $\varepsilon(\alpha_1,\alpha_2)$.

\textbf{Bar degree 3:}
\[
d([e^{\alpha_1} \mid e^{-\alpha_1-\alpha_2} \mid e^{\alpha_2}]
\otimes \eta_{12} \wedge \eta_{23})
= [e^0] = [1]
\]
since both adjacent pairs have inner product $-1$. This realizes the
Weyl group relation in the bar complex.

The nontrivial differentials encode the $A_2$ root system structure:
simple poles occur precisely for pairs of roots whose sum is again a root.
\end{computation}

\begin{proposition}[$D_4$ bar complex and triality; \ClaimStatusProvedHere]
\label{prop:lattice:bar-D4}
For the $D_4$ root lattice, the bar complex
$\barBgeom(\Vlat_{D_4})$ admits an action of the triality group
$S_3 \cong \operatorname{Out}(D_4)$ permuting the three conjugacy classes of
roots. The induced action on homology
$H_*(\barBgeom(\Vlat_{D_4}))$ reflects the triality symmetry of the
$\AffKM{so}_8$ current algebra at level~1.
\end{proposition}

\begin{proof}
The $D_4$ root lattice has 24 roots, all of the form $\pm e_i \pm e_j$ ($i \neq j$).  These form a single orbit under the Weyl group $W(D_4)$.  The outer automorphism group $\operatorname{Out}(D_4) \cong S_3$ acts as graph automorphisms of the $D_4$ Dynkin diagram, permuting the three ``legs'' and hence permuting the three nontrivial cosets of $D_4^*/D_4 \cong (\mathbb{Z}/2\mathbb{Z})^2$ (vector, spinor, and conjugate spinor representations).  Note: the spinor and conjugate spinor vectors lie in $D_4^* \setminus D_4$, not in the root lattice itself.

\textbf{Step 1: Triality action on bar complex.}
The triality automorphism $\sigma$ acts on the lattice as a graph automorphism of the $D_4$ Dynkin diagram, permuting the three ``legs.''  This extends to an automorphism of $\Vlat_{D_4}$ via $\sigma(e^\alpha) = e^{\sigma(\alpha)}$.  Since the bar construction is functorial (any algebra automorphism induces a coalgebra automorphism of the bar complex), $\sigma$ acts on $\barBgeom(\Vlat_{D_4})$.

\textbf{Step 2: Action on homology.}
In bar degree~1, the generators are $[e^\alpha]$ for each root $\alpha$.  The triality permutes these among the three classes of 8.  In bar degree~2, the differential $d([e^{\alpha_1} \mid e^{\alpha_2}]) = \varepsilon(\alpha_1,\alpha_2)\,[e^{\alpha_1+\alpha_2}]$ (when $\langle\alpha_1,\alpha_2\rangle = -1$) intertwines with $\sigma$ since $\sigma$ preserves inner products and the cocycle is equivariant (the standard Frenkel--Lepowsky--Meurman cocycle for $D_4$ is triality-invariant up to coboundary).

\textbf{Step 3: Conclusion.}
The $S_3$-action on $H_*(\barBgeom(\Vlat_{D_4}))$ is the representation-theoretic manifestation of the triality symmetry: the three fundamental representations of $\mathfrak{so}_8$ (vector, spinor, conjugate spinor) are permuted by outer automorphisms, and this permutation lifts to the bar complex level.
\end{proof}

\begin{proposition}[$E_8$ bar complex and self-duality; \ClaimStatusProvedHere]
\label{prop:lattice:bar-E8}
For the $E_8$ root lattice, the bar complex $\barBgeom(\Vlat_{E_8})$ is
self-dual under Koszul duality (Theorem~\ref{thm:lattice:koszul-dual}),
reflecting the unimodularity and self-duality of $E_8$.
The bar complex decomposes as
\[
\barBgeom(\Vlat_{E_8})
= \bigoplus_{\gamma \in E_8} \barBgeom(\Vlat_{E_8})_\gamma
\]
with $|E_8/2E_8| = 2^8 = 256$ sectors modulo~$2E_8$. The 240 roots of
$E_8$ each contribute a 1-dimensional summand in bar degree~1.
\end{proposition}

\begin{proof}
\textbf{Step 1: Lattice decomposition.}
The bar complex decomposes along the $E_8$-grading: $\barBgeom(\Vlat_{E_8}) = \bigoplus_{\gamma \in E_8} \barBgeom(\Vlat_{E_8})_\gamma$.  The zero sector $\gamma = 0$ is the bar complex of the rank-8 Heisenberg algebra $\barBgeom(\cH_8)$, whose bar differential vanishes (the Heisenberg OPE has only a double pole, so $\cH_8^! \simeq \mathrm{Sym}^{ch}(V^*)$ by Theorem~\ref{thm:heisenberg-koszul-dual-early}).

\textbf{Step 2: Root contributions.}
In bar degree~1, each root $\alpha$ of $E_8$ contributes a generator $[e^\alpha]$ in the sector $\gamma = \alpha$.  Since $E_8$ has 240 roots, $\dim \barBgeom^1 = 240$ (in the root sectors).  In bar degree~2, the differential is $d([e^{\alpha_1} \mid e^{\alpha_2}]) = \varepsilon(\alpha_1,\alpha_2)\,[e^{\alpha_1+\alpha_2}]$ when $\langle\alpha_1,\alpha_2\rangle = -1$.  For $E_8$, each root has 56 neighbors at inner product $-1$ (from the $E_8$ root system geometry: the 126 pairs of roots summing to a root correspond to 56 neighbors per root, counted with multiplicity).

\textbf{Step 3: Self-duality.}
By Theorem~\ref{thm:lattice:koszul-dual}, $(\Vlat_{E_8}^\varepsilon)^! = (\Vlat_{E_8}^{\varepsilon^{-1}})^c$.  Since $E_8$ is unimodular ($E_8 = E_8^*$), the cocycle $\varepsilon$ for $E_8$ satisfies $\varepsilon^{-1} \sim \varepsilon$ (cohomologous): since $\mathbb{C}^\times$ is $2$-divisible, the bimultiplicative cocycle $\varepsilon$ has a square root $\varepsilon^{1/2}$, and the coboundary $\delta(\varepsilon^{1/2})$ witnesses $\varepsilon^{-1} \sim \varepsilon$.  Therefore $(\Vlat_{E_8})^! \cong \Vlat_{E_8}$ as a coalgebra.

\textbf{Step 4: Sector count.}
Modulo $2E_8$, the lattice has $|E_8/2E_8| = 2^{\operatorname{rank}(E_8)} = 2^8 = 256$ classes (since $\Lambda/2\Lambda \cong (\mathbb{Z}/2\mathbb{Z})^d$ for any rank-$d$ lattice), giving 256 sectors in the bar complex.
\end{proof}

\begin{theorem}[Unimodular lattice self-duality; \ClaimStatusProvedHere]\label{thm:lattice:unimodular-self-dual}
Let $\Lambda$ be an even unimodular lattice.  Then $\Vlat_\Lambda$ is Koszul self-dual:
\[
\Vlat_\Lambda^! \;\cong\; \Vlat_\Lambda
\]
and the bar-cobar quasi-isomorphism $\Omega^{\mathrm{ch}}(\barBgeom(\Vlat_\Lambda)) \xrightarrow{\sim} \Vlat_\Lambda$ witnesses this self-duality.
\end{theorem}

\begin{proof}
By Theorem~\ref{thm:lattice:koszul-dual}, $\Vlat_\Lambda^! = (\Vlat_\Lambda^{\varepsilon^{-1}})^c$.  For $\Lambda$ unimodular, the cocycles $\varepsilon$ and $\varepsilon^{-1}$ are cohomologous: since $\C^\times$ is $2$-divisible, every bimultiplicative $2$-cocycle $\varepsilon : \Lambda \times \Lambda \to \C^\times$ has a square root $\varepsilon^{1/2}$, and the coboundary $\delta(\varepsilon^{1/2})$ witnesses $\varepsilon^{-1} \sim \varepsilon$.  Hence $\Vlat_\Lambda^{\varepsilon^{-1}} \cong \Vlat_\Lambda^\varepsilon$ as vertex algebras.  The coalgebra structure on the dual side matches the algebra structure on the original side (since the cocycle is self-inverse up to coboundary), giving $\Vlat_\Lambda^! \cong \Vlat_\Lambda$.

The bar-cobar QI follows from Theorem~\ref{thm:lattice:koszul-morphism}: the twisting morphism $\tau$ is Koszul, so $\Omega^{\mathrm{ch}}(\barBgeom(\Vlat_\Lambda)) \xrightarrow{\sim} \Vlat_\Lambda$.
\end{proof}


% ======================================================================
% Section 6: Koszul Duality for Lattice Algebras
% ======================================================================

\section{Koszul Duality for Lattice Algebras}
\label{sec:lattice:koszul}

\begin{definition}[Inverse cocycle]\label{def:lattice:inverse-cocycle}
For a cocycle $\varepsilon \colon \Lambda \times \Lambda \to \C^\times$,
the \textbf{inverse cocycle} is
$\varepsilon^{-1}(\alpha,\beta) := \varepsilon(\alpha,\beta)^{-1}$.
The cocycle condition is preserved under inversion, and the commutator
satisfies
$c_{\varepsilon^{-1}}(\alpha,\beta) = c_\varepsilon(\alpha,\beta)^{-1}$.
In particular, if $\varepsilon$ is symmetric, so is $\varepsilon^{-1}$.
\end{definition}

\begin{theorem}[Koszul dual of lattice vertex algebra; \ClaimStatusProvedHere]
\label{thm:lattice:koszul-dual}
The Koszul dual of the lattice $\Eone$-chiral algebra
$\Vlat_\Lambda^\varepsilon$ is:
\begin{equation}\label{eq:lattice:koszul-dual}
(\Vlat_\Lambda^\varepsilon)^!
\;=\; (\Vlat_\Lambda^{\varepsilon^{-1}})^c,
\end{equation}
the coalgebra structure on the lattice vertex algebra with inverse
cocycle. More precisely:
\begin{enumerate}[label=(\roman*)]
\item The underlying graded vector space of
  $(\Vlat_\Lambda^\varepsilon)^!$ is $\Vlat_\Lambda^*$.
\item The coalgebra structure is determined by $\varepsilon^{-1}$:
  the coproduct $\Delta \colon \Vlat_\Lambda^* \to
  \Vlat_\Lambda^* \otimes \Vlat_\Lambda^*$ has structure constants
  $\varepsilon(\alpha,\beta)^{-1}$.
\item For symmetric $\varepsilon$, the Koszul dual is again a
  (co)lattice vertex algebra with the same symmetry type.
\end{enumerate}
\end{theorem}

\begin{proof}
The Koszul dual is computed via the bar-cobar adjunction. The
product $\mu \colon \Vlat_\Lambda \otimes \Vlat_\Lambda \to \Vlat_\Lambda$
has structure constants $\varepsilon(\alpha,\beta)$ from the group
algebra multiplication $e^\alpha \cdot e^\beta =
\varepsilon(\alpha,\beta)\,e^{\alpha+\beta}$.
The Koszul dual coalgebra is extracted from the bar complex
$\bar{B}^{\mathrm{ch}}(\Vlat_\Lambda^\varepsilon)$.  In the bar construction,
each tensor factor is desuspended by $s^{-1}$, and the bar differential
involves the product $\mu$ composed with Koszul signs from the
desuspension.  For the lattice algebra, where $\mu(e^\alpha \otimes
e^\beta) = \varepsilon(\alpha,\beta)\,e^{\alpha+\beta}$ with
$|\varepsilon(\alpha,\beta)| = \pm 1$, the desuspension introduces a
sign $(-1)^{|\alpha||\beta|}$ that, combined with the bilinearity of
$\varepsilon$, produces the cocycle inversion:
the bar-complex coproduct has structure constants
$\varepsilon(\alpha,\beta)^{-1}$.
Since $\varepsilon^{-1}$ is again a valid cocycle,
$(\Vlat_\Lambda^{\varepsilon^{-1}})^c$ is a valid coalgebra, and the
Koszul duality exchanges the two.
\end{proof}

\begin{theorem}[Koszul morphism for lattice algebras; \ClaimStatusProvedHere]
\label{thm:lattice:koszul-morphism}
The canonical twisting morphism
$\tau \colon \barBgeom(\Vlat_\Lambda^\varepsilon) \to
\Vlat_\Lambda^\varepsilon$
is a Koszul morphism: the twisted complex
\[
\barBgeom(\Vlat_\Lambda^\varepsilon)
\otimes_\tau \Vlat_\Lambda^\varepsilon
\]
is acyclic.
\end{theorem}

\begin{proof}
Filter by lattice weight. The associated graded is the bar complex of
the Heisenberg algebra (forgetting the lattice part), which is acyclic
by the Heisenberg bar-cobar quasi-isomorphism
(cf.\ the free field computations in Chapter~\ref{ch:lattice}).
The spectral sequence associated to the lattice weight filtration
converges (the lattice $\Lambda$ is positive-definite, so the
filtration is bounded below in each total degree) and degenerates at
$E_2$. The cocycle $\varepsilon$ contributes phases but does not affect
the acyclicity argument, since the underlying vector spaces are the
same for all cocycles.
\end{proof}

\begin{remark}[Self-duality for unimodular lattices]
\label{rem:lattice:self-duality}
When $\Lambda$ is unimodular ($\Lambda = \Lambda^*$) and the cocycle
$\varepsilon$ satisfies $\varepsilon^{-1} \sim \varepsilon$ (cohomologous),
the Koszul dual $(\Vlat_\Lambda^\varepsilon)^!$ is isomorphic to
$\Vlat_\Lambda^\varepsilon$ itself (as a coalgebra). This is the case
for $E_8$: the lattice vertex algebra $\Vlat_{E_8}$ is
\emph{Koszul self-dual}. (Note: the Heisenberg algebra $\mathcal{H}_\kappa$ is
Koszul but \emph{not} self-dual: $\mathcal{H}_\kappa^! \simeq \mathrm{Sym}^{ch}(V^*)$
by Theorem~\ref{thm:heisenberg-koszul-dual-early}. Self-duality is a special
feature of the $E_8$ lattice structure.)
\end{remark}

\begin{table}[ht]
\centering
\caption{Master table: vertex algebra families as lattice constructions.}
\label{tab:lattice:master-table}
\begin{tabular}{l l l l}
\textbf{Family} & \textbf{Lattice origin} & \textbf{Operation} & \textbf{Chapter} \\
\hline
Heisenberg $\cH_d$
  & Any $\Lambda$, rank $d$ & Zero sector & \ref{ch:lattice} \\
$\AffKM{g}$ at level 1
  & Root lattice $\Lambda_\fg$ & FKS (Thm~\ref{thm:lattice:frenkel-kac}) & 9 \\
$\AffKM{g}$ at level $k$
  & Rescaled lattice & Sublattice embedding & 9 \\
$W$-algebras $\WAlg{k}{g}$
  & $\Lambda_\fg$ & BRST/DS reduction & 10 \\
$bc$ ghosts
  & $\Z\alpha$, $\langle\alpha,\alpha\rangle = -1$ & Odd lattice & 8 \\
Lattice VOA $\Vlat_\Lambda$
  & $\Lambda$ (even) & Direct construction & \ref{ch:lattice} \\
Orbifold $(\Vlat_\Lambda)^\sigma$
  & $\Lambda/\sigma$ & Fixed-point subalgebra & \ref{sec:lattice:functorial} \\
Overlattice $\Vlat_{\Lambda_I}$
  & $\Lambda_I \supset \Lambda$ & Extension by $I$ & \ref{sec:lattice:functorial}
\end{tabular}
\end{table}


% ======================================================================
% Section 7: Functorial Properties and Lattice Operations
% ======================================================================

\section{Functorial Properties and Lattice Operations}
\label{sec:lattice:functorial}

This section is the conceptual heart of the ``engine'' vision.
Each lattice operation from
Definition~\ref{def:lattice:operations} lifts to a functorial operation
on vertex algebras, and these lifts are compatible with the bar
construction.

\subsection{Direct sum}

\begin{theorem}[Tensor product from direct sum; \ClaimStatusProvedHere]
\label{thm:lattice:direct-sum}
For even lattices $\Lambda_1, \Lambda_2$ with cocycles
$\varepsilon_1, \varepsilon_2$:
\begin{equation}\label{eq:lattice:tensor-product}
\Vlat_{\Lambda_1 \oplus \Lambda_2}^{\varepsilon_1 \otimes \varepsilon_2}
\;\cong\;
\Vlat_{\Lambda_1}^{\varepsilon_1}
\otimes \Vlat_{\Lambda_2}^{\varepsilon_2}
\end{equation}
where $(\varepsilon_1 \otimes \varepsilon_2)
((\alpha_1,\alpha_2),(\beta_1,\beta_2)) =
\varepsilon_1(\alpha_1,\beta_1)\,\varepsilon_2(\alpha_2,\beta_2)$.
\end{theorem}

\begin{proof}
The state space decomposes as
$\cH_{\Lambda_1 \oplus \Lambda_2} \otimes \C[\Lambda_1 \oplus \Lambda_2]
\cong (\cH_{\Lambda_1} \otimes \C[\Lambda_1]) \otimes
(\cH_{\Lambda_2} \otimes \C[\Lambda_2])$.
The vertex operator for $e^{(\alpha_1,\alpha_2)}$ factors as
$Y(e^{\alpha_1}, z) \otimes Y(e^{\alpha_2}, z)$
since the oscillators for $\Lambda_1$ and $\Lambda_2$ commute
(orthogonality of the direct sum).
\end{proof}

\begin{corollary}[K\"unneth for bar complexes; \ClaimStatusProvedElsewhere]\label{cor:lattice:kunneth}
Under the isomorphism~\eqref{eq:lattice:tensor-product}:
\[
\barBgeom(\Vlat_{\Lambda_1} \otimes \Vlat_{\Lambda_2})
\simeq
\barBgeom(\Vlat_{\Lambda_1})
\otimes \barBgeom(\Vlat_{\Lambda_2}).
\]
\end{corollary}

\begin{proof}
The bar construction of a tensor product of vertex algebras is
quasi-isomorphic to the tensor product of bar constructions, by
the Eilenberg--Zilber/K\"unneth theorem for the bar construction
(cf.\ the algebraic foundations in Chapter~2). The lattice grading
ensures the relevant spectral sequence converges.
\end{proof}

\subsection{Sublattice embedding}

\begin{proposition}[Sublattice maps; \ClaimStatusProvedHere]\label{prop:lattice:sublattice}
An embedding of even lattices $\iota \colon \Lambda' \hookrightarrow \Lambda$
induces:
\begin{enumerate}[label=(\roman*)]
\item An embedding of vertex algebras
  $\iota_* \colon \Vlat_{\Lambda'}^{\iota^*\varepsilon}
  \hookrightarrow \Vlat_\Lambda^\varepsilon$
  (restriction of the cocycle).
\item A map of bar complexes
  $\barBgeom(\iota_*) \colon
  \barBgeom(\Vlat_{\Lambda'}) \to \barBgeom(\Vlat_\Lambda)$.
\item On homology, the induced map respects the lattice grading:
  the $\gamma$-sector of $H_*(\barBgeom(\Vlat_{\Lambda'}))$ maps to
  the $\iota(\gamma)$-sector of $H_*(\barBgeom(\Vlat_\Lambda))$.
\end{enumerate}
\end{proposition}

\begin{proof}
(i) The inclusion $\iota$ induces $\C[\Lambda'] \hookrightarrow \C[\Lambda]$
and $\cH_{\Lambda'} \hookrightarrow \cH_\Lambda$. The cocycle on
$\Lambda'$ is the restriction $\iota^*\varepsilon$.

(ii) The bar construction is functorial in the algebra.

(iii) The lattice grading is preserved by $\iota$.
\end{proof}

\subsection{Overlattice extension}

\begin{theorem}[Overlattice vertex algebra; \ClaimStatusProvedElsewhere]\label{thm:lattice:overlattice}
Let $\Lambda$ be an even lattice and $I \subset D(\Lambda) = \Lambda^*/\Lambda$
an isotropic subgroup. The overlattice $\Lambda_I$ is even, and:
\begin{equation}\label{eq:lattice:overlattice}
\Vlat_{\Lambda_I}
\;\cong\; \bigoplus_{\gamma \in I} \Vlat_\Lambda^\gamma
\end{equation}
where $\Vlat_\Lambda^\gamma$ denotes the $\gamma$-twisted module
for $\Vlat_\Lambda$ (the sector of states with lattice charge in the
coset $\gamma + \Lambda$).
\end{theorem}

\begin{proof}[Proof sketch]
The state space of $\Vlat_{\Lambda_I}$ decomposes under the sublattice
$\Lambda \subset \Lambda_I$ as
$\Vlat_{\Lambda_I} = \bigoplus_{\gamma \in \Lambda_I/\Lambda}
\cH \otimes \C[\gamma + \Lambda]$.
Each summand $\cH \otimes \C[\gamma + \Lambda]$ is the twisted module
$\Vlat_\Lambda^\gamma$. The isotropicity of $I$ ensures that the
extended multiplication is well-defined (the cocycle extends from
$\Lambda$ to $\Lambda_I$) and that the vertex algebra axioms hold.
\end{proof}

\begin{example}[Heterotic lattices]\label{ex:lattice:heterotic}
The two even unimodular lattices of rank~16 arise as overlattices:
\begin{enumerate}[label=(\roman*)]
\item $E_8 \oplus E_8$: already unimodular, no extension needed.
\item $D_{16}^+ = (D_{16})_I$ where $I = \{0, [\text{spinor}]\} \cong \Z/2\Z$
  is the isotropic subgroup generated by the spinor class in
  $D(\Lambda_{D_{16}}) \cong (\Z/2\Z)^2$.
\end{enumerate}
The corresponding vertex algebras $\Vlat_{E_8 \oplus E_8}$ and
$\Vlat_{D_{16}^+}$ give the two heterotic string theories. Their bar
complexes are non-isomorphic despite having the same rank and
determinant.
\end{example}

\subsection{Orbifold construction}

\begin{construction}[Orbifold vertex algebra]
\label{constr:lattice:orbifold}
For $\sigma \in \operatorname{Aut}(\Lambda)$ of finite order $N$,
the \textbf{orbifold vertex algebra} is the fixed-point subalgebra:
\[
(\Vlat_\Lambda^\varepsilon)^\sigma
= \{v \in \Vlat_\Lambda^\varepsilon : \sigma \cdot v = v\},
\]
where $\sigma$ acts on $\Vlat_\Lambda$ by
$\sigma(e^\alpha) = e^{\sigma(\alpha)}$ and $\sigma(\alpha_{-n}) =
(\sigma\alpha)_{-n}$ on Heisenberg oscillators.

When $\sigma$ preserves the cocycle ($\varepsilon(\sigma\alpha,\sigma\beta)
= \varepsilon(\alpha,\beta)$), this is a vertex subalgebra.
\end{construction}

\begin{table}[ht]
\centering
\caption{Functorial operations on lattice algebras and their bar complexes.}
\label{tab:lattice:bar-operations}
\begin{tabular}{l l l l}
\textbf{Operation} & \textbf{Vertex algebra} & \textbf{Bar complex} & \textbf{Koszul dual} \\
\hline
Direct sum $\Lambda_1 \oplus \Lambda_2$
  & $\Vlat_{\Lambda_1} \otimes \Vlat_{\Lambda_2}$
  & $\barBgeom_1 \otimes \barBgeom_2$
  & $\Vlat_{\Lambda_1}^! \otimes \Vlat_{\Lambda_2}^!$ \\[3pt]
Sublattice $\Lambda' \hookrightarrow \Lambda$
  & $\Vlat_{\Lambda'} \hookrightarrow \Vlat_\Lambda$
  & $\barBgeom_{\Lambda'} \to \barBgeom_\Lambda$
  & $\Vlat_\Lambda^! \twoheadrightarrow \Vlat_{\Lambda'}^!$ \\[3pt]
Overlattice $\Lambda_I \supset \Lambda$
  & $\Vlat_{\Lambda_I} = \bigoplus_{I} \Vlat_\Lambda^\gamma$
  & $\barBgeom_{\Lambda_I} = \bigoplus_{I} \barBgeom_\gamma$
  & $\Vlat_{\Lambda_I}^! = \bigoplus_{I} (\Vlat_\Lambda^\gamma)^!$ \\[3pt]
Orbifold $\Vlat_\Lambda^\sigma$
  & $(\Vlat_\Lambda)^\sigma$
  & $\barBgeom(\Vlat_\Lambda)^\sigma$
  & $(\Vlat_\Lambda^!)^\sigma$ \\[3pt]
Rescaling $\sqrt{k}\Lambda$
  & $\Vlat_{\sqrt{k}\Lambda}$
  & Level-$k$ bar complex
  & Level-$k$ dual \\[3pt]
BRST reduction
  & $H^0_{Q_{\mathrm{DS}}}(\Vlat_\Lambda)$
  & $H^*_{Q_{\mathrm{DS}}}(\barBgeom(\Vlat_\Lambda))$
  & W-algebra Koszul dual
\end{tabular}
\end{table}

\begin{remark}[The derivation chain]\label{rem:lattice:derivation-chain}
The four operations --- direct sum, sublattice embedding, overlattice
extension, and orbifold --- generate the following derivation chain that
produces every example in Part~2:
\[
\text{Lattice } \Lambda
\xrightarrow{\text{FKS}}
\text{Kac--Moody}
\xrightarrow{\text{BRST}}
\text{$W$-algebras}
\xrightarrow{\text{orbifold}}
\text{non-simply-laced}
\]
More precisely:
\begin{enumerate}[label=(\arabic*)]
\item Start with an even positive-definite lattice $\Lambda$.
\item The Frenkel--Kac--Segal theorem (Theorem~\ref{thm:lattice:frenkel-kac})
  gives the level-1 affine algebra $L(\AffKM{g}, 1)$.
\item Higher levels arise from sublattice embeddings
  $\sqrt{k}\,\Lambda \hookrightarrow \Lambda$ (rescaling the bilinear
  form by~$k$).
\item $W$-algebras arise by quantum Drinfeld--Sokolov reduction (BRST
  cohomology of the affine algebra with respect to a nilpotent
  subalgebra).
\item Non-simply-laced algebras arise from orbifolds of simply-laced
  lattice constructions.
\item Tensor products of these give further examples (direct sum of
  lattices).
\end{enumerate}
At each step, the bar complex transforms functorially, so the Koszul
duality theory of Part~1 applies throughout.
\end{remark}


% ======================================================================
% Section 8: Chiral Hochschild Cohomology
% ======================================================================

\section{Chiral Hochschild Cohomology of Lattice Algebras}
\label{sec:lattice:hochschild}

\begin{theorem}[Lattice Hochschild cohomology; \ClaimStatusProvedHere]\label{thm:lattice:hochschild}
For an even integral lattice $\Lambda$:
\begin{equation}\label{eq:lattice:hochschild}
HH^n_{\mathrm{ch}}(\Vlat_\Lambda, \Vlat_\Lambda)
\;\cong\; H^n(\Lambda, \Vlat_\Lambda^{\Lambda})
\end{equation}
where $\Vlat_\Lambda^\Lambda$ denotes the $\Lambda$-invariants (the
degree-zero sector under the lattice grading) and $H^n(\Lambda, -)$ is
group cohomology of the free abelian group $\Lambda$.
\end{theorem}

\begin{proof}
The lattice vertex algebra decomposes as
$\Vlat_\Lambda = \bigoplus_{\alpha \in \Lambda} \Vlat_\Lambda^\alpha$
where $\Vlat_\Lambda^\alpha = \cH \otimes e^\alpha$ is the
$\alpha$-isotypic component.

The chiral Hochschild complex respects this grading:
\[
CC^n_{\mathrm{ch}}(\Vlat_\Lambda, \Vlat_\Lambda)
= \bigoplus_{\alpha_1, \ldots, \alpha_n, \beta}
\Hom(\Vlat_\Lambda^{\alpha_1} \otimes \cdots
\otimes \Vlat_\Lambda^{\alpha_n}, \Vlat_\Lambda^\beta).
\]
The differential preserves $\sum_i \alpha_i = \beta$, and the
resulting spectral sequence (filtering by lattice degree)
collapses at $E_2$ to give:
\[
HH^*_{\mathrm{ch}}(\Vlat_\Lambda, \Vlat_\Lambda)
= \bigoplus_{\alpha \in \Lambda}
HH^*_{\mathrm{ch}}(\Vlat_\Lambda, \Vlat_\Lambda^\alpha).
\]
For $\alpha = 0$, this identifies with the group cohomology
$H^*(\Lambda, \Vlat_\Lambda^{\Lambda})$ via the standard bar resolution
for $\Lambda \cong \Z^d$.
\end{proof}

\begin{corollary}[Unimodular case; \ClaimStatusProvedHere]\label{cor:lattice:hochschild-unimodular}
For a unimodular lattice $\Lambda$ (so $D(\Lambda) = 0$), the chiral
Hochschild cohomology reduces to:
\[
HH^n_{\mathrm{ch}}(\Vlat_\Lambda, \Vlat_\Lambda)
\cong H^n(\Lambda, \cH)
\]
where $\cH$ is the Heisenberg Fock space (the $\Lambda$-invariant
sector when $\Lambda = \Lambda^*$).
\end{corollary}

\begin{example}[$A_1$ Hochschild cohomology]\label{ex:lattice:hochschild-A1}
For $\Lambda = A_1 \cong \Z$, the group cohomology of $\Z$ is:
\[
H^n(\Z, M) =
\begin{cases}
M^\Z & n = 0, \\
M_\Z & n = 1, \\
0 & n \geq 2.
\end{cases}
\]
Hence $HH^n_{\mathrm{ch}}(\Vlat_{A_1}, \Vlat_{A_1})$ is concentrated
in degrees 0 and~1. This reflects the rigidity of the $\AffKM{sl}_2$
current algebra at level~1: it has a 1-dimensional space of first-order
deformations (the level) and no higher obstructions.
\end{example}


% ======================================================================
% Section 9: Higher Genus
% ======================================================================

\section{Higher Genus}
\label{sec:lattice:higher-genus}

At higher genus, the lattice grading on the bar complex persists and
interacts with the geometry of moduli spaces to produce modular objects.

\begin{proposition}[Genus-1 partition function; \ClaimStatusProvedHere]
\label{prop:lattice:genus-1}
For an even lattice $\Lambda$ of rank $d$, the genus-1 partition
function of $\Vlat_\Lambda$ is:
\begin{equation}\label{eq:lattice:genus-1-Z}
Z_1(\Vlat_\Lambda, \tau)
= \frac{\Theta_\Lambda(\tau)}{[\eta(\tau)]^d}
\end{equation}
where $\Theta_\Lambda(\tau) = \sum_{\alpha \in \Lambda}
q^{\langle\alpha,\alpha\rangle/2}$ is the theta series of $\Lambda$
(with $q = e^{2\pi i \tau}$) and $\eta(\tau) = q^{1/24}
\prod_{n=1}^\infty (1 - q^n)$ is the Dedekind eta function.
\end{proposition}

\begin{proof}
The partition function traces over the state space:
$Z_1 = \operatorname{tr}_{\Vlat_\Lambda} q^{L_0 - c/24}$.
The Heisenberg sector contributes $1/\eta(\tau)^d$ (the bosonic
oscillator partition function for $d$ free bosons), and the lattice
sum $\sum_{\alpha \in \Lambda} q^{h_{e^\alpha}}
= \sum_\alpha q^{\langle\alpha,\alpha\rangle/2} = \Theta_\Lambda(\tau)$
accounts for the lattice momentum modes.
\end{proof}

\begin{theorem}[Modular invariance; \ClaimStatusProvedHere]\label{thm:lattice:modular-invariance}
For an even unimodular lattice $\Lambda$ of rank $d$, the partition
function $Z_1(\Vlat_\Lambda, \tau)$ is $S$-invariant:
$Z_1(\Vlat_\Lambda, -1/\tau) = Z_1(\Vlat_\Lambda, \tau)$.
Under $T: \tau \mapsto \tau+1$, the partition function transforms as
$Z_1(\tau+1) = e^{-\pi i d/12}\, Z_1(\tau)$; full
$\mathrm{SL}_2(\Z)$-invariance holds if and only if $d \equiv 0 \pmod{24}$.
\end{theorem}

\begin{proof}
By Poisson summation, $\Theta_\Lambda(-1/\tau) =
(\tau/i)^{d/2} |\Lambda^*/\Lambda|^{1/2}\, \Theta_{\Lambda^*}(\tau)$.
For unimodular $\Lambda = \Lambda^*$, this gives
$\Theta_\Lambda(-1/\tau) = (\tau/i)^{d/2}\, \Theta_\Lambda(\tau)$.
Combined with $\eta(-1/\tau) = \sqrt{\tau/i}\,\eta(\tau)$,
we obtain $Z_1(-1/\tau) = Z_1(\tau)$.

For the $T$-transformation $\tau \mapsto \tau + 1$:
$\Theta_\Lambda(\tau+1) = \Theta_\Lambda(\tau)$ (since
$\langle\alpha,\alpha\rangle \in 2\Z$ for even $\Lambda$) and
$\eta(\tau+1) = e^{\pi i/12}\eta(\tau)$, giving
$Z_1(\tau+1) = e^{-\pi i d/12} Z_1(\tau)$. For $d \equiv 0 \pmod{24}$
this is exact invariance; in general, $Z_1$ transforms as a modular
function of the appropriate level.

For $d = 8$ ($E_8$): the $T$-phase is $e^{-8\pi i/12} = e^{-2\pi i/3} \neq 1$, so
$Z_1(\tau+1) = e^{-2\pi i/3}\, Z_1(\tau)$: the partition function is $S$-invariant
but \emph{not} $T$-invariant.  Equivalently,
$Z_1 = E_4(\tau)/\eta(\tau)^8 = j(\tau)^{1/3}$ transforms under $T$ with a
cube-root-of-unity multiplier, making it a modular function for the congruence
subgroup~$\Gamma(3)$ but not for the full~$\mathrm{SL}_2(\Z)$.
Full modular invariance holds only for $d \equiv 0 \pmod{24}$
(e.g., the Leech lattice with $d = 24$).
\end{proof}

\begin{remark}[Higher genus and Siegel modular forms]
\label{rem:lattice:higher-genus}
At genus $g$, the partition function generalizes to
\[
Z_g(\Vlat_\Lambda, \Omega)
= \frac{\Theta_\Lambda^{(g)}(\Omega)}{[\det(\operatorname{Im}\Omega)]^{d/2}}
\cdot (\text{oscillator contribution})
\]
where $\Omega$ is the period matrix, $\Theta_\Lambda^{(g)}$ is the
Siegel theta series, and the denominator involves the determinant of the
imaginary part of $\Omega$. For unimodular even lattices,
$\Theta_\Lambda^{(g)}$ transforms as a Siegel modular form of weight
$d/2$ for $\operatorname{Sp}(2g, \Z)$, as discussed in
Chapter~14.

The key point for this monograph: the lattice grading on the bar complex
persists at all genera. At genus $g$, the bar differential acquires
corrections from the cohomology of $\overline{\cM}_{g,n}$, but these
corrections preserve the total lattice degree. The spectral sequence
computing bar complex homology at genus $g$ has $E_1$ page given by
the genus-0 bar complex (with its lattice decomposition), and the
higher differentials $d_r$ for $r \geq 1$ encode the genus-$g$
quantum corrections studied in Part~1.
\end{remark}

\section{Catalog of Lattice Koszul Dual Pairs}\label{sec:lattice:catalog}

The lattice engine produces Koszul dual pairs for every even lattice.
Table~\ref{tab:lattice:koszul-catalog} summarizes the results.

\begin{table}[ht]
\centering
\caption{Lattice Koszul dual pairs and genus-1 bar complex data.}
\label{tab:lattice:koszul-catalog}
\begin{tabular}{l c c c c}
\textbf{Lattice $\Lambda$} & \textbf{Rank} & \textbf{$\Vlat_\Lambda^!$} & \textbf{$\Theta_\Lambda(\tau)$} & \textbf{Self-dual?} \\
\hline
$A_1$ & 1 & $\Vlat_{A_1^*}$ & $\theta_3(2\tau)$ & No \\
$A_2$ & 2 & $\Vlat_{A_2^*}$ & $\theta_3^{A_2}(\tau)$ & No \\
$D_4$ & 4 & $\Vlat_{D_4^*}$ & $\frac{1}{2}(\theta_3^4+\theta_4^4)$ & No \\
$D_8$ & 8 & $\Vlat_{D_8^*}$ & $\Theta_{D_8}(\tau)$ & No \\
$E_6$ & 6 & $\Vlat_{E_6^*}$ & $\Theta_{E_6}(\tau)$ & No \\
$E_7$ & 7 & $\Vlat_{E_7^*}$ & $\Theta_{E_7}(\tau)$ & No \\
$E_8$ & 8 & $\Vlat_{E_8}$ & $E_4(\tau)$ & Yes \\
$E_8 \oplus E_8$ & 16 & $\Vlat_{E_8 \oplus E_8}$ & $E_4(\tau)^2$ & Yes \\
$D_{16}^+$ & 16 & $\Vlat_{D_{16}^+}$ & $E_4(\tau)^2$ & Yes \\
$\mathrm{Leech}$ & 24 & $\Vlat_{\mathrm{Leech}}$ & $\Theta_{\mathrm{Leech}} = j-744$ & Yes \\
\end{tabular}
\end{table}

\begin{proposition}[Koszul self-duality criterion; \ClaimStatusProvedHere]\label{prop:lattice:self-dual-criterion}
$\Vlat_\Lambda$ is Koszul self-dual if and only if $\Lambda$ is unimodular (i.e., $\Lambda = \Lambda^*$) and the FLM cocycle $\varepsilon$ is self-inverse up to coboundary.  The even unimodular lattices in rank $d$ exist only for $d \equiv 0 \pmod{8}$; the count is $1$ for $d = 8$, $2$ for $d = 16$, $24$ for $d = 24$ (Niemeier lattices), and grows rapidly thereafter.
\end{proposition}

\begin{proof}
The ``only if'' direction: if $\Vlat_\Lambda^! \cong \Vlat_\Lambda$, then by Theorem~\ref{thm:lattice:koszul-dual}, $\Vlat_{\Lambda}^{\varepsilon^{-1}} \cong \Vlat_\Lambda^\varepsilon$, which requires $\Lambda^* = \Lambda$ (unimodularity, since the dual lattice enters via the coalgebra structure) and $\varepsilon^{-1} \sim \varepsilon$.  The ``if'' direction is Theorem~\ref{thm:lattice:unimodular-self-dual}.  The existence condition $d \equiv 0 \pmod 8$ for even unimodular lattices is classical (signature constraints from the Milnor--Husemoller classification).
\end{proof}

\begin{proposition}[$D_4$ and triality; \ClaimStatusProvedHere]\label{prop:lattice:D4-triality}
The $D_4$ lattice is \emph{not} unimodular: $\det(\mathrm{Gram}_{D_4}) = 4$ and $D_4^*/D_4 \cong (\Z/2\Z)^2$.  In particular, $\Vlat_{D_4}$ is \emph{not} Koszul self-dual; its Koszul dual is $\Vlat_{D_4^*}$.

The overlattice $D_4^+ = D_4 \cup (D_4 + s)$ (adding the spinor class $s = \tfrac{1}{2}(1,1,1,1)$) is unimodular but \emph{odd}: $\langle s, s\rangle = 1 \notin 2\mathbb{Z}$, so $D_4^+ \cong \mathbb{Z}^4$ as lattices.  Even unimodular overlattices $D_n^+$ exist only for $n \equiv 0 \pmod{8}$ (e.g., $D_8^+ \cong E_8$).  Triality (the $S_3$ symmetry of the $D_4$ Dynkin diagram) permutes the three nontrivial cosets in $D_4^*/D_4$ --- vector, spinor, and conjugate spinor --- and acts as an outer automorphism of the $D_4$ root system.
\end{proposition}

\begin{proof}
The discriminant computation $D_4^*/D_4 \cong (\Z/2\Z)^2$ shows $D_4 \neq D_4^*$, so Theorem~\ref{thm:lattice:koszul-dual} gives $\Vlat_{D_4}^! \cong \Vlat_{D_4^*} \not\cong \Vlat_{D_4}$.  The overlattice $D_4^+$ is unimodular ($[D_4^+ : D_4] = 2$ and $\det = 4/4 = 1$) but odd ($\langle s, s\rangle = 1$), so Theorem~\ref{thm:lattice:unimodular-self-dual} does not apply.  The triality action on $D_4^*/D_4$ is the standard $S_3$-symmetry permuting the three nontrivial cosets.
\end{proof}

\begin{remark}[Distinguishing $E_8 \oplus E_8$ from $D_{16}^+$]
The bar complexes of $\Vlat_{E_8 \oplus E_8}$ and $\Vlat_{D_{16}^+}$ are both self-dual (both are even unimodular lattices of rank~16), but they are \emph{non-isomorphic} as coalgebras.  At genus~1, the two lattices are \emph{indistinguishable}: both have theta function $\Theta_\Lambda(\tau) = E_4(\tau)^2$, since the space~$M_8(\mathrm{SL}_2(\Z))$ is one-dimensional.  The distinction first appears at genus~2: the Siegel theta functions $\Theta^{(2)}_{E_8 \oplus E_8}$ and $\Theta^{(2)}_{D_{16}^+}$ are distinct Siegel modular forms of weight~8, providing a complete invariant distinguishing the two heterotic lattices.  This illustrates how the genus-$g$ bar complex captures strictly finer information than the genus-1 bar complex.
\end{remark}
